\documentclass[../textbook.tex]{subfiles}
\begin{document}
\bookchapter{模型评估}{ch:evaluation-statistics}

评估要回答的是一个有条件的决策问题：在什么任务与用户分布上，某个系统是否比基线更有效，代价和风险是否可以接受，以及现有证据还有多大不确定性。有限测试集上的分数只是这个问题的一部分。任务定义错误、标签失真或测试污染都可能产生很精确却不对应目标的数字。

本章将指标、测量误差和统计推断放在同一条证据链上。检索与智能体的组件指标分别见\bookxref{ch:information-retrieval}、\bookxref{ch:rag-systems}及\bookxref{ch:agent-runtime}；逐例观测通过指标与统计推断形成系统比较，人工和模型评审用于内容判定，在线随机实验用于估计干预效果。

\section{评估对象}

\subsection{总体、协议及目标量}
\term{待估目标量}{Estimand}{}是统计分析真正要估计的对象。设任务输入及其参考信息为 $X\sim\mathcal P$，固定系统为 $f_v$，随机生成条件为 $\omega$，评估协议为 $e$，逐例效用为 $u$。目标可以写为
\begin{equation}
 \mu_v=\mathbb E_{X\sim\mathcal P,\omega}\left[u(X,f_v(X;\omega);e)\right].
 \label{eq:ch41-estimand}
\end{equation}
这里的系统版本 $v$ 不只是模型名称，还包括提示、分词器、解码、检索、工具和输出解析等配置。改变协议 $e$ 也可能改变目标量，例如把不适用项计零与从分母排除，会得到不同含义的分数。

\term{基准测试}{Benchmark}{}是在约定任务和协议下形成的可比较坐标。部署中的任务效用还包含错误损失、完成时间、重试、人工介入与资源成本；基准得分不能不经论证就替代式\eqref{eq:ch41-estimand}。即使某个题库中所有答案都正确，也不能据此推断另一语言、工具环境或时间范围的任务都正确。

\begin{table}[htbp]
\centering
\caption{主要统计对象。}\label{tab:evaluation-statistics-symbols}
\begin{tabularx}{\linewidth}{@{}lX@{}}
\toprule
符号 & 含义\\
\midrule
$n,G,m_g$ & 记录数、独立组数、第 $g$ 组记录数\\
$s_i^{(A)},s_i^{(B)},d_i$ & 同一案例两系统分数与配对差\\
$p,\hat p$ & 总体通过概率与样本通过比例\\
$\alpha,z_{1-\frac{\alpha}{2}}$ & 显著性水平与标准正态分位数\\
$\Delta,\delta$ & 系统期望差与预先规定的最小效应或非劣界\\
$\mathcal P,\mathcal Q$ & 目标任务分布与实际抽样分布\\
\bottomrule
\end{tabularx}
\end{table}

\subsection{能力及证据分层}
知识评估需要限定知识时间、领域与可访问来源；语言评估关注理解、表达和语言覆盖；推理评估需要检验中间约束与最终结论；长上下文评估需要控制证据位置、干扰与组合跨度；工具使用评估还需要判定参数、授权、执行结果和副作用。通用能力、领域能力与安全对齐必须按任务目的分别定义，不能用大量普通问答平均掉少数高损失失败。

组件指标解释局部机制，端到端任务指标观察整体完成情况，鲁棒性和安全切片观察特定失败面，在线随机实验估计真实流量中的干预效果。评估方案应分别覆盖候选召回、回答忠实度、工具权限与最终任务完成。

\begin{figure}[htbp]
\centering
\begin{tikzpicture}[>=stealth,every node/.style={font=\small},b/.style={draw,rounded corners,align=center,text width=29mm,minimum height=10mm}]
\node[b] (a) at (0,0) {目标总体与决策\\明确效用和损失};
\node[b] (b) at (3.8,0) {冻结案例与系统\\保留逐例输出};
\node[b] (c) at (7.6,0) {测量与评审\\校准误差};
\node[b] (d) at (7.6,-1.8) {统计单位与区间\\配对、聚类};
\node[b] (e) at (3.8,-1.8) {切片与约束\\质量、风险、成本};
\node[b] (f) at (0,-1.8) {决策与受控实验\\晋级、保持、回退};
\draw[->] (a)--(b);\draw[->] (b)--(c);\draw[->] (c)--(d);\draw[->] (d)--(e);\draw[->] (e)--(f);
\end{tikzpicture}
\caption{从任务定义到决策证据。每一层都保留自己的适用条件。}\label{fig:ch41-evidence}
\end{figure}

\section{常用数据集及基准协议}

\label{sec:ch41-benchmarks}
\subsection{环境构建}
评测材料至少有三种形态。静态数据集提供问题、参考答案或标签；交互基准还提供工具、初始状态、时间限制和执行后的验收器；综合指数则汇总多个基准的结果。三类材料分别组织静态问题、交互过程与跨基准汇总。本节选取研究与模型发布中有代表性的基准，选型时应核对其任务分布。

对静态选择题，核心观测通常是答案是否正确。对交互任务，\term{评测运行框架}{Evaluation Harness}{}负责组织提示、工具调用、环境恢复、预算与结果收集。同一个模型换用不同运行框架，可能执行不同的搜索与修复策略。因此，一个可解释的分数必须绑定
\begin{equation}
 \mathcal E=(D_v,M_v,H_v,B,R,A),
 \label{eq:ch41-benchmark-protocol}
\end{equation}
其中 $D_v$ 为数据或环境版本，$M_v$ 为模型版本，$H_v$ 为运行框架，$B$ 为词元、时间、工具调用和尝试次数预算，$R$ 为逐例计分规则，$A$ 为聚合规则。安全控制及回退模型也属于系统配置。两个结果只要这些要素存在实质差异，就不能仅凭同一个基准名称推断为受控对照。

\subsection{知识、科学及数学}
\term{研究生级难检索问答}{Graduate-Level Google-Proof Question Answering}{GPQA}由生物、物理和化学领域专家编写多项选择题，考察专门知识与推理；Diamond 是其中的特定子集，不能与完整版本混写。主要指标是答案准确率。题目难以通过普通检索解决，是基准设计目标，不构成永远不存在网络答案的保证\cite{rein2024gpqa}。

\term{人类最后的考试}{Humanity's Last Exam}{HLE}覆盖多学科的高难度封闭式问题，包含多模态材料，公开定稿集为2500题，并有单独保留的私有测试材料。报告中需要区分文本子集、完整输入、工具可用性及数据修订版本。其准确率与置信度校准考察不同性质：准确率衡量答案正确性，校准衡量置信度与观察频率的对应；开放式科研能力另按交互任务评估\cite{hle2026benchmark}。

FrontierMath 面向高难度数学问题。Tier 4 是特定难度层，不等于全部题目；v2 修订过题目与集合，因此版本必须随分数保留。官方协议允许使用计算工具，并提交可验证答案。答案验证检查数值或符号对象，证明审查逐步检查逻辑有效性\cite{epoch2026frontiermathv2}。

\subsubsection{基础知识及中文评估}
\term{大规模多任务语言理解}{Massive Multitask Language Understanding}{MMLU}使用跨学科选择题，为知识与推理提供基础坐标；MMLU-Pro 强化了推理要求并调整题目与选项，属于另一套基准。比较时须同时固定少样本示例、选项顺序、答案解析和学科权重\cite{hendrycks2021mmlu,wang2024mmlupro}。如果模型已经接近某个题集的上限，少量分数变化提供的区分信息会减少；仍可将其用于回归检测，但选型还需要更贴近目标任务的材料。

C-Eval 与\term{中文大规模多任务语言理解}{Chinese Massive Multitask Language Understanding}{CMMLU}分别覆盖52个学科与67个主题，适合补充中文及本地知识的评估。C-Eval Hard、验证集与测试集也必须分别标识。中文长文写作、引用问答与业务流程使用各自的任务评估\cite{huang2023ceval,li2023cmmlu}。

\subsubsection{基础数学及竞赛题}
GSM8K 以多步基础算术应用题为主；MATH 包含竞赛数学题及参考解答。通常依据解析后的最终答案计算正确率，但数值格式、符号等价判定和工具使用都属于协议，不能仅报告模型名称\cite{cobbe2021gsm8k,hendrycks2021math}。MATH-500 与完整测试集的样本构成不同，报告中应保留子集名称。

\term{美国数学邀请赛}{American Invitational Mathematics Examination}{AIME}是逐年举行的竞赛，每套15题，答案为0至999的整数；它不是一个自带统一模型采样协议的静态数据集。模型评测应写明年份、I卷或II卷、是否合并、每题尝试次数及答案选择方法\cite{maa2026aime}。例如一套15题中多答对一题就改变约6.67个百分点，故小题集的年度差异与抽样变异尤其不能忽略。

\subsubsection{函数生成及时间分组的代码题}
HumanEval 根据函数签名和文档字符串要求生成 Python 实现，通过测试估计 pass@$k$；其组合估计见式\eqref{eq:ch41-passk}。LiveCodeBench 持续汇集竞赛编程题，并区分代码生成、自修复等场景；用于比较时应固定发布版本、题目日期窗口和测试裁剪方式\cite{chen2021codeeval,jain2024livecodebench}。日期划分用于降低时间泄漏风险，还需结合模型训练来源检查。一次生成与参考答案引导的重试分别记录预算和结果。

\subsubsection{多模态、长上下文及指令遵循}
\term{大规模多学科多模态理解}{Massive Multi-discipline Multimodal Understanding}{MMMU}将大学学科问题与图表等视觉材料结合，按答案正确性计分。需要固定图像分辨率、多图处理及测试划分，避免把视觉预处理差异误认为纯文本推理差异\cite{yue2024mmmu}。MMMU 与 MMMU-Pro 分属不同协议。公开题目与标签的状态也会改变：不能因早期论文采用隐藏测试标签，就永久假定当前测试答案仍未公开。

LongBench 原版混合长文问答、摘要、检索与代码等任务，分别采用 F1、ROUGE 或相似度等指标；LongBench v2 改用更强调长上下文推理的多项选择题。原版总分不是统一的答案准确率，v2 也不能被当作原题集仅增加长度\cite{bai2024longbench,bai2024longbench2}。对长度的计量还应注明字符、词还是具体分词器的词元。

\term{多轮共指消解}{Multi-Round Coreference Resolution}{MRCR}要求在含相近请求的长对话中找回指定的既有回答。OpenAI 公开版本按针数及长度范围分组，计分器先检查指定前缀，再计算文本相似度；它不是一般问答的二元准确率\cite{openai2025mrcr}。这类合成检索成绩需要与多文档推理互补，不能据此声称整个长上下文都被可靠理解。

\term{指令遵循评估}{Instruction-Following Evaluation}{IFEval}利用程序检查词数、关键词与输出形式等可验证约束。它同时区分提示级和指令级、严格和宽松两组口径：提示级要求一条提示内所有约束均通过，指令级按单个约束汇总；宽松计分只容许协议列出的有限输出变换，并不是任意人工放宽要求\cite{zhou2023ifeval}。格式、事实准确性与业务价值分别评价。



\subsection{代码、浏览及计算机操作}
SWE-bench 根据真实软件仓库的问题描述要求系统生成修补，并在仓库环境中运行测试；Verified 是经过人工筛选的500个实例子集。主要指标为解决比例。该结果受基础提交、依赖环境、测试验收及智能体运行框架影响；函数生成与仓库修复分别评价，结果按实际验收测试的范围解释\cite{swebench2026verified}。

Terminal-Bench 系列通过终端环境中的可执行任务评估智能体。需要把通用终端任务与科学工作流分支分开：例如发布材料中的 Terminal-Bench 4.0 与 Terminal-Bench Science 0.1 是不同版本和任务分布，不能合并成一项“代码准确率”。科学工作流还涉及数据分析、模拟与模型拟合；最终产物是否满足验收器，比回答是否声称完成更有意义\cite{openai2026astralaunch}。

BrowseComp 以难以定位、但答案较短且可核验的信息检索问题为对象，原始集合包含1266题。它适合研究多步搜索与证据定位，主要报告答案准确率；浏览器、搜索源、访问时间和预算都会影响结果。长篇报告的论证质量、引用覆盖和用户需求理解仍需另外评估\cite{openai2025browsecomp}。

OSWorld 2.0 是跨软件的长程计算机操作基准，可按完整集或离线子集评价，也区分部分奖励与二元完成判定。ScreenSpot-Pro 则侧重屏幕元素定位；定位与工作流完成具有不同的验收对象。前者需要检查最终状态和多步约束，后者的测量对象更接近动作选择中的视觉定位组件\cite{osworld2026v2,openai2026astralaunch}。

\subsection{交互推理及计分单位}
ARC-AGI-3 要求智能体探索陌生的交互环境并从反馈中学习，计分包含相对于人类的行动效率。它与 ARC-AGI-1、ARC-AGI-2 的静态任务形式不同；三个版本使用各自的任务与计分规则，结果按对应环境解释\cite{arc2026agi3}。

即使分数都被写为百分数，也必须检查分子和分母。设任务 $i$ 有 $m_i$ 个验收项，权重 $w_{ij}\geq0$ 且 $\sum_jw_{ij}=1$，验收结果为 $r_{ij}\in\{0,1\}$。可以构造任务部分得分与全通过指标
\begin{equation}
 r_i^{\mathrm{partial}}=\sum_{j=1}^{m_i}w_{ij}r_{ij},\qquad
 r_i^{\mathrm{complete}}=\prod_{j=1}^{m_i}r_{ij}.
 \label{eq:ch41-partial-complete}
\end{equation}
这里假定每个验收项均为完成任务所必需，且没有额外一票否决规则。两个任务各有五项等权要求，分别满足四项与两项，则平均部分得分为60\%，完整任务成功率却为0。部分得分与完整任务成功率反映不同的验收要求。

\subsection{模型发布评测}
\phantomsection\label{q:ch41:example:01}

\label{sec:ch41-release-benchmarks}
以下记录2026年9月8日核实的 GPT-6 Astra 与 Claude Fable 5.1 发布材料，用于说明如何阅读模型发布表。表中数值采用厂商发布时的报告结果。表\ref{tab:evaluation-release-scores}选取共同或具有代表性的项目；空项表示本表没有对应发布数值，不能填成零。

\begin{longtable}{@{}p{.38\linewidth}p{.16\linewidth}p{.16\linewidth}p{\dimexpr.30\linewidth-6\tabcolsep\relax}@{}}
\caption{发布时的代表性评测结果}\label{tab:evaluation-release-scores}\\
\toprule 基准或计分项 & Astra & Fable 5.1 & 指标及限定\\\midrule
\endfirsthead
\toprule 基准或计分项 & Astra & Fable 5.1 & 指标及限定\\\midrule
\endhead
Terminal-Bench 4.0 & 57.9\% & 55.8\% & 终端任务得分；运行框架须另核对\\
Terminal-Bench Science 0.1 & 64.6\% & 52.6\% & 科学工作流得分\\
HLE，有工具 & 57.2\% & 65.0\% & 答案准确率；工具与裁判协议不同\\
HLE，无工具 & — & 60.9\% & 与有工具配置分开\\
AutomationBench & 41.4\% & 31.4\% & 业务流程任务通过比例\\
GDPval-AA v2 & — & 1853 & Elo，非百分数\\
CursorBench 3.2.0 & — & 73.4\% & Cursor 编码任务评分\\
OSWorld 2.0，Astra 发布协议 & 72.6\% & — & v2026.08.08 离线集，部分得分\\
OSWorld 2.0，Fable 发布协议 & — & 77.9\% & 8月任务版及修复，部分得分\\
同上，严格完成口径 & — & 41.7\% & 全部要求满足的任务比例\\
GPQA Diamond & 96.0\% & — & 科学选择题准确率\\
FrontierMath Tier 4 v2 & 97.6\% & — & 指定数学题集得分\\
ARC-AGI-3 & 99.9\% & — & 交互推理与行动效率\\
BrowseComp & 91.5\% & — & 浏览检索答案准确率\\
\bottomrule
\end{longtable}
数据来源分别为 OpenAI 的 Astra 发布表与 Anthropic 的 Fable 5.1 发布主表\cite{openai2026astralaunch,anthropic2026fable51launch}。表中采用 Fable 发布主表的七个基准、九个计分行，以及 Astra 发布表的代表项目。Astra 引言把 FrontierMath 成绩取整为98\%，这里保留正式表中的97.6\%。



GPT-6 Astra 的发布表还包含专业工作、科学健康、安全和长上下文等评测：例如 Agents' Last Exam、BenchCAD、DeepSWE v1.1、FrontierCode 1.1、HealthBench Professional、GeneBench Pro、LifeSciBench、MRCR v2，以及内部设计、数据科学和数据库迁移任务。Artificial Analysis 的 Intelligence Index 与 Coding Agent Index 是综合指数，不是单个数据集。内部任务结果按厂商提供的题目范围与协议解释\cite{openai2026astralaunch}。

读发布表时，尤其要保留下列限定。第一，Astra 表中的 OSWorld 2.0 明确为 v2026.08.08 离线集合与部分得分，不能与不同任务修改及评分设置的结果直接相减。第二，Astra 发布表报告各推理档位中的最高分；它不是固定成本下的统一对照。推理强度的模型相关语义与对照方法见\bookxref{ch:verifiable-reasoning}。第三，安全干预后改由另一个模型完成的任务，应被解释为包含回退路径的系统结果，不能归给一个未经切换的模型。\footnote{Astra 发布材料还注明，部分标在 Fable 名下的 ScreenSpot-Pro、ExploitGym 数值来自安全限制不同的 Mythos 配置。模型名称、访问配置与实际执行者必须共同记录。}\cite{openai2026astralaunch}

Fable 5.1 发布案例中的 GDPval-AA v2 根据职业任务产物的盲配对比较形成 Elo；AutomationBench 检查跨应用业务流程的最终状态；CursorBench 3.2.0 使用 Cursor 的编码任务与运行框架。Elo 是相对评级，1853不能读成1853\%的准确率，跨日期的评级还可能因参评系统与拟合结果变化而改变。它们分别提供产物偏好、流程验收和编码工作流证据，不应直接平均为一个百分数\cite{anthropic2026fable51card}。

Fable 的系统卡还说明了影响解释的实验条件：HLE 分别运行有工具与无工具配置，使用固定裁判，并对检索到题库答案的轨迹重新判错；OSWorld 被安全机制中止的指定任务计零，其他受干预的特定任务可能由较早模型接替。Terminal-Bench Science 0.1 各模型报告的标准误约为3.5至4.5个百分点，主要不确定性来自任务差异。这是标准误范围，不是每个模型统一的95\%置信区间\cite{anthropic2026fable51card}。同一道任务的多次运行也不能都当成新的独立任务，本章后文的聚类与重复测量方法仍然适用。

比较设计由实际决策目标确定。若目标是选择固定月度预算下的代码助手，就应在相同任务和成本约束下比较质量；若目标是研究最大可达能力，可以提高预算，但需同时报告尝试次数、运行时间与成本。跨厂商榜单通常不能提供同题配对结果，因此不应直接套用本章配对检验给分数差赋予显著性。

\begin{algorithm}[选择并冻结评测组合]
\textbf{输入：}目标任务分布、错误损失、候选系统、时间与费用预算。\par
\textbf{输出：}版本化基准组合、逐例验收规则与统计计划。
\begin{enumerate}
 \item 将目标拆成知识、推理、检索、工具动作和最终产物等可观察能力，确定哪些约束必须同时满足。
 \item 选取相关公开基准，并加入来自真实业务、与训练和开发过程隔离的保留案例；记录每组材料的来源与权重。
 \item 冻结式\eqref{eq:ch41-benchmark-protocol}中的全部要素，特别是工具、预算、失败计分、安全控制与回退路径。
 \item 在开发材料上校验验收器和答案解析；冻结后保留所有测试案例，包括超时、拒答、无效产物与环境故障的分类记录。
 \item 保存逐例观测，按预先定义的独立单位计算区间及切片结果；同时报告质量、延迟、费用和高损失失败。
\end{enumerate}
\textbf{停止与边界：}达到预定样本及预算后停止。公开分数无法提供的配置或逐例信息应明确缺失，不通过猜测补齐。评测组合变化后重新定义目标量。
\end{algorithm}


\section{分类、概率及决策代价}

\subsection{混淆矩阵及平均口径}
\term{混淆矩阵}{Confusion Matrix}{}以真实标签为行、预测标签为列。二分类中的真阳性、假阳性、真阴性和假阴性分别记为 TP、FP、TN、FN。\term{准确率}{Accuracy}{}、\term{精确率}{Precision}{}和\term{召回率}{Recall}{}为
\begin{equation}
 \mathrm{Acc}=\frac{\mathrm{TP}+\mathrm{TN}}n,\quad
 P=\frac{\mathrm{TP}}{\mathrm{TP}+\mathrm{FP}},\quad
 R=\frac{\mathrm{TP}}{\mathrm{TP}+\mathrm{FN}}.
\end{equation}
\term{F1 分数}{F1 Score}{}是精确率和召回率的调和平均：
\begin{equation}
 F_1=\frac{2PR}{P+R}=\frac{2\mathrm{TP}}{2\mathrm{TP}+\mathrm{FP}+\mathrm{FN}}.
 \label{eq:ch41-f1}
\end{equation}
分母为零时须显式规定未定义、排除或计零，并报告受影响类别。把“本次没有这种任务”与“这种任务全失败”都记成零会改变结论。

设样本数 $n=100$：实际正类20条，其中12条预测为正；实际负类80条，其中4条预测为正。因此 $\mathrm{TP}=12,\mathrm{FN}=8,\mathrm{FP}=4,\mathrm{TN}=76$。准确率为 $0.88$，正类精确率 $\frac{12}{16}=0.75$，召回率 $\frac{12}{20}=0.60$，F1 为 $\frac{24}{36}=\frac{2}{3}$。若把负类也按一对其余类别计算，负类 F1 为 $\frac{152}{164}\approx0.9268$。

\term{宏平均}{Macro Average}{}使各类别等权，本例 Macro-F1 约为 $\frac{0.6667+0.9268}{2}=0.7967$。\term{微平均}{Micro Average}{}先汇总各类计数，再计算指标。在单标签、类别全集齐全的多分类中，每次错误同时贡献一个 FP 和一个 FN，故 Micro-F1 等于准确率，本例为 $0.88$。这一等式不应推广到多标签或只选择部分类别的评估。

\subsection{阈值及概率校准}
若 $p=\Pr(Y=1\mid x)$ 是已校准条件概率，误报和漏报损失为 $c_{\mathrm{FP}},c_{\mathrm{FN}}>0$，预测正类的条件风险是 $c_{\mathrm{FP}}(1-p)$，预测负类的风险是 $c_{\mathrm{FN}}p$。选择前者当且仅当
\begin{equation}
 c_{\mathrm{FP}}(1-p)<c_{\mathrm{FN}}p
 \quad\Longleftrightarrow\quad
 p>\frac{c_{\mathrm{FP}}}{c_{\mathrm{FP}}+c_{\mathrm{FN}}}.
 \label{eq:ch41-threshold}
\end{equation}
例如误报损失为9、漏报为1，阈值为0.9。这个推导假设正确决策损失为零且损失不随个体变化；有额外行动、拒答或个体成本时，应比较全部行动的条件风险。

\term{概率校准}{Probability Calibration}{}要求预测为某概率水平的样本具有相应事件频率。排序好不等于校准好。一个模型把所有概率从 $p$ 单调变换为 $p^2$，排序可以保持，式\eqref{eq:ch41-threshold}的概率含义却不再自动成立。二分类 Brier 分数为 $n^{-1}\sum_i(p_i-y_i)^2$；若真实概率为 $q$，其条件期望等于 $(p-q)^2+q(1-q)$，因此在 $p=q$ 处最小。

\term{期望校准误差}{Expected Calibration Error}{ECE}常将预测分桶，计算各桶预测均值与事件比例差的加权绝对值。它依赖分桶规则，粗分桶可能掩盖局部错误，样本稀少时又有较大估计噪声。阈值和校准器都应在开发或校准数据上确定，最终测试数据用于评估冻结方案。

\subsection{排序曲线}
提高阈值使预测正类集合缩小，因此召回率非增；精确率不具有一般的单调保证，因为删除的下一个样本可能是真阳性，也可能是假阳性。\term{受试者工作特征曲线}{Receiver Operating Characteristic}{ROC}描绘真阳性率与假阳性率；其面积可解释为随机正例分数高于随机负例的概率，平分时通常计半分。

\term{精确率—召回率曲线}{Precision--Recall Curve}{PR}尤其需要同时报告正类基率。按排序位置 $j$ 记精确率 $P_j$ 和召回增量 $R_j-R_{j-1}$，一种\term{平均精确率}{Average Precision}{AP}定义为 $\sum_jP_j(R_j-R_{j-1})$，不等于任意插值方式的曲线面积。曲线概括阈值族，部署阈值仍应由实际损失与约束决定。

\section{生成指标的测量对象}

\subsection{概率及词面一致性}


\term{困惑度}{Perplexity}{PPL}将逐词元平均负对数似然指数化：
\begin{equation}
 \mathrm{PPL}=\exp\left(-\frac1T\sum_{t\in\mathcal S}\log p(x_t\mid x_{<t})\right),\qquad T=|\mathcal S|.
\end{equation}
$\mathcal S$ 是实际计分位置集合。分词器、文本、上下文窗口、模板、起止符和掩码共同规定“每词元”的单位；重叠滑窗不能无意重复计分。跨分词器的 PPL 不能直接按数值大小比较模型能力。其概率模型基础见\bookxref{ch:rev-autoregressive}。

\term{完全匹配}{Exact Match}{EM}比较规范化后的答案是否相等。规范化必须声明大小写、空白、标点、数值单位及多参考规则。\term{词元 F1}{Token F1}{}令候选和参考的词元多重集合重合数为 $m$，候选与参考长度分别为 $c,r$，则 $P=\frac{m}{c},R=\frac{m}{r},F_1=\frac{2m}{c+r}$。重复词元的匹配次数不能超过参考中的出现次数。

\term{双语评估替代指标}{Bilingual Evaluation Understudy}{BLEU}使用经参考计数截断的 $n$ 元组精确率 $p_n$，并施加长度惩罚\cite{papineni2002bleu}：
\begin{equation}
 \mathrm{BLEU}=\mathrm{BP}\exp\left(\sum_{n=1}^{N_g}w_n\log p_n\right),\quad
 \mathrm{BP}=\begin{cases}1,&c>r,\\e^{1-\frac{r}{c}},&0<c\leq r.\end{cases}
\end{equation}
$w_n\geq0$ 且总和为1，$c,r$ 为按协议得到的候选和有效参考长度。通常需在语料层累加计数，而非随意平均逐句 BLEU。\footnote{BLEU 的名称来自机器翻译评估。将其用于其他生成任务时，仍需论证参考文本及词面重合与实际任务质量的关系，不能仅沿用一个熟悉的指标名称。}零匹配时的平滑、空输出和多参考长度选择都应固定。

\term{面向召回的摘要评估指标}{Recall-Oriented Understudy for Gisting Evaluation}{ROUGE}是指标族\cite{lin2004rouge}。ROUGE-$n$ 的基本召回形式为匹配 $n$ 元组数除以参考 $n$ 元组数；ROUGE-L 则利用最长公共子序列，具体输出召回、精确率或 F 分数仍须声明。词面相近可以对应事实相反，否定词、数值和单位错误也可能只改变很少词元。

\begin{example}[词面指标的手工计算]
\label{q:ch41:example:03}
以参考词元序列“甲 乙 丙 丁”、候选“甲 乙 戊”为例，$m=2,c=3,r=4$，Token F1 为 $\frac{4}{7}$，ROUGE-1 召回率为 $\frac{1}{2}$，EM 为零。只计算一元与二元 BLEU、各权重 $\frac{1}{2}$ 且不平滑时，$p_1=\frac{2}{3},p_2=\frac{1}{2}$，故分数为 $e^{-\frac{1}{3}}\sqrt{\frac{1}{3}}\approx0.4137$。这些指标描述词元匹配，事实错误由内容判定评价。
\end{example}

\subsection{事实、执行及拒答}
事实评估应把回答分解为可判定论断，分别标记证据支持、冲突与未知。引用存在、来源相关和来源支持论断是三个不同事件。引用精确率的分母可以是已给出的引用，论断覆盖率的分母则是需要证据的论断；二者不得混用。检索覆盖与答案忠实性的关系见\bookxref{ch:information-retrieval}及\bookxref{ch:rag-systems}。

对于数学、代码和结构化操作，\term{可执行验证器}{Executable Verifier}{}以可执行约束或测试判定结果。测试结果按实际输入与约束的覆盖范围解释。评估执行环境需隔离待测代码和副作用，超时、异常及无效输出必须按预定规则进入分母，而不是只保留成功返回的候选。

若对一题独立同分布生成 $n$ 个候选，其中 $c$ 个通过固定验证器，从这些候选均匀无放回选 $k\leq n$ 个，全部失败的子集数为 $\binom{n-c}{k}$，总子集数为 $\binom nk$，故
\begin{equation}
 \widehat{\mathrm{pass@}k}=1-\frac{\binom{n-c}{k}}{\binom nk}.
 \label{eq:ch41-passk}
\end{equation}
在同分布独立采样前提下，这个子集平均对 $k$ 次独立尝试至少一次通过的概率无偏\cite{chen2021codeeval}。例如 $n=5,c=2,k=2$，结果为 $1-\frac{3}{10}=0.7$；直接代入样本通过比例得到 $1-(\frac{3}{5})^2=0.64$，两种估计并不相同。相关采样、不同温度混合或自适应重试会改变目标分布。pass@$k$评价候选覆盖，最终交付成功率还取决于选择器。

拒答评估还应区分不可回答任务上的正确弃权，与可回答任务上的过度拒绝。安全失败、事实错误和风格偏好应分维度报告，不能以更流畅的文风抵消一个禁止行动。

\section{评审员测量模型}

\subsection{人工量表及模型裁判}
\term{评分量表}{Rubric}{}需要将“好回答”拆成可观察维度，为各等级提供锚点、正反例及不适用规则。人工评审先独立评分再仲裁，保留原始分歧；只保存仲裁标签会隐藏量表歧义。盲评隐藏候选身份、运行成本等无关信息，候选顺序随机化，必要时限制高风险材料的访问与保留。

\term{模型评审}{LLM-as-a-Judge}{}可扩展开放式判断，但会受位置、长度、风格、自我偏好和提示注入影响\cite{zheng2023judge}。交换顺序后必须先映射回原候选身份，再比较判定。相同事实的长短改写可以识别冗长偏好；跨模型家族的交叉设计用于识别自我偏好。一个总体人工一致率无法替代这些专门对照。

对人工判定不通过的样本，假通过率定义为 $\Pr(J=1\mid H=0)$；在裁判判定通过的样本中，错误比例则为 $\Pr(H=0\mid J=1)$，两者分母不同。裁判模型、量表、语言或候选输出分布改变后，应重新估计这两个量。只有两个顺序都通过才接受，是一种保守决策规则，可能降低假通过，也可能提高假拒绝；两个判断通常相关，不能把两次错误率直接相乘。

\subsection{一致性及胜率计算}


\term{科恩一致性系数}{Cohen's Kappa}{$\kappa$}以边际分布为基准校正观测一致率\cite{cohen1960agreement}。若观测一致率为 $p_o$，两位评审在类别 $c$ 的边际比例为 $p_{1c},p_{2c}$，则
\begin{equation}
 p_e=\sum_cp_{1c}p_{2c},\qquad \kappa=\frac{p_o-p_e}{1-p_e},\quad p_e<1.
\end{equation}
\begin{example}[二分类评审的一致性]
\label{q:ch41:example:02}
例如100条二分类评分中，两人都判通过40条、都判不通过30条，甲通过乙不通过20条，甲不通过乙通过10条。于是 $p_o=0.7$，甲通过率0.6、乙通过率0.5，$p_e=0.6\times0.5+0.4\times0.5=0.5$，$\kappa=0.4$。该指标衡量超过偶然基准的一致程度；类别极不均衡会改变这个基准。序数、多评审和缺失评分需要适合其结构的统计量。
\end{example}

成对偏好中若胜、负、平分别为 $W,L,T$，半分平局胜率为 $\frac{W+\frac{T}{2}}{W+L+T}$。例如 $(12,6,2)$ 得到 $0.65$。单条贡献可以是 $0,\frac{1}{2},1$，因此它是有界均值而非普通伯努利比例，不能直接把“13胜20题”当作真实二项计数套用所有二项区间。

\section{区间估计}

\subsection{Wilson 区间推导}
\begin{assumption}[独立二项比例]
本节假设每个统计单位产生一次通过或失败，单位相互独立、通过概率相同，样本来自目标抽样机制，系统与判定规则已冻结。多轮会话、重复用户和多个随机种子并不自动满足这些条件。
\end{assumption}

令 $K\sim\operatorname{Binomial}(n,p)$，$\hat p=\frac{K}{n}$，则 $\mathbb E\hat p=p$、$\operatorname{Var}(\hat p)=\frac{p(1-p)}{n}$。直接把 $p$ 换成 $\hat p$ 得到 Wald 区间，在边界处可能退化或越界。\term{威尔逊区间}{Wilson Interval}{}通过反解关于候选 $p$ 的分数检验不等式获得\cite{wilson1927inference}：
\begin{align}
 (\hat p-p)^2&\leq \frac{z^2p(1-p)}{n},\\
 (n+z^2)p^2-(2n\hat p+z^2)p+n\hat p^2&\leq0.
\end{align}
取二次方程两根之间的区间，得到
\begin{equation}
 \frac{\hat p+\frac{z^2}{2n}\,\pm\ z\sqrt{\frac{\hat p(1-\hat p)}{n}+\frac{z^2}{4n^2}}}{1+\frac{z^2}{n}},
 \qquad z=z_{1-\frac{\alpha}{2}}.
 \label{eq:ch41-wilson}
\end{equation}
Wilson 区间采用正态分位数近似，实际覆盖率随 $n,p$ 变化。\term{置信区间}{Confidence Interval}{CI}的频率学派解释是：在相同抽样机制反复取样并按同一规则构造区间时，长期覆盖固定真值的比例接近声明水平。它不表示观察到本区间后，固定参数以95\%的概率随机落在其中。

对于 $n=20,k=20,z=1.96$，上界为1，下界为 $\frac{20}{20+3.8416}\approx0.8389$。二十次全通过仍允许相当大的未观测失败率。另从零失败的精确二项关系 $(1-q)^n=\alpha$，得到失败率 $q$ 的单侧上界 $1-\alpha^{\frac{1}{n}}$；当 $q$ 小时近似为 $-\frac{\log\alpha}{n}$，95\%时约为 $\frac{3}{n}$。单侧界与前面的双侧 Wilson 界使用不同构造，不能混为一个指标。

区间仅反映给定模型下的抽样误差。判定器偏差、样本代表性和训练污染不会因为 $n$ 很大就自动消失。

\subsection{配对差及效应大小}
同一批案例比较 A 与 B，令 $d_i=s_i^{(A)}-s_i^{(B)}$。在独立案例和有限方差条件下，
\begin{equation}
 \bar d=\frac1n\sum_i d_i,\quad
 s_d^2=\frac1{n-1}\sum_i(d_i-\bar d)^2,\quad
 \widehat{\operatorname{SE}}(\bar d)=\frac{s_d}{\sqrt n.}
 \label{eq:ch41-paired}
\end{equation}
其优势来自 $\operatorname{Var}(A-B)=\operatorname{Var}(A)+\operatorname{Var}(B)-2\operatorname{Cov}(A,B)$。两模型通常同时在难题上失败，正相关可以降低配对差的方差；把结果当作无关样本会丢弃这部分信息。

例如100个独立案例中，A独胜15条、B独胜5条，其余80条一致。差值总和10、平方和20，所以 $\bar d=0.10$，$s_d^2=\frac{20-100\times0.1^2}{99}=\frac{19}{99}$，标准误约0.0438，正态近似95\%区间约为 $[0.014,0.186]$。该算例支持正向平均差，但区间下界未达到假设的最小实用提升0.02，因而“区别于零”和“达到实用标准”是不同结论。

\term{置换检验}{Permutation Test}{}在零假设下交换配对中的系统身份，并比较统计量。独立配对且零假设赋予交换对称性时，可枚举或抽样翻转 $d_i$ 符号；仅假设均值为零而不具备相应对称或随机化条件，并不足以自动保证任意符号检验有效。$p$ 值是零假设及检验模型下至少同样极端结果的概率，不是零假设为真的概率。

\subsection{聚类、Bootstrap 及多种子}
\term{自助重采样}{Bootstrap}{}从经验样本重复有放回抽样，以近似统计量的抽样分布\cite{efron1979bootstrap}。计算非线性指标如 Macro-F1 时，应在每次重采样后重算指标，而非只重采样一个已经聚合的数字。

\term{簇重采样}{Cluster Bootstrap}{}将用户、文档等独立组作为抽样单位，保留组内记录共同出现。若每组有 $m$ 条记录，单条方差 $\sigma^2$，组内任意两条相关系数为 $\rho$、组间独立，则 $n=Gm$ 条记录均值的方差为
\begin{equation}
 \operatorname{Var}(\bar X)=\frac{1}{n^2}G\left[m\sigma^2+m(m-1)\rho\sigma^2\right]
 =\frac{\sigma^2}{n}[1+(m-1)\rho].
 \label{eq:ch41-cluster}
\end{equation}
方括号是这一简化模型的设计效应。例如10位用户各10条请求、$\rho=0.5$，设计效应为5.5，按方差等价得到有效样本量约 $\frac{100}{5.5}=18.2$，不能视为100个独立用户。

组大小不一致时还必须明确估计请求平均还是用户平均。二者分别为
\begin{equation}
 \bar s_{\mathrm{request}}=\frac{\sum_g\sum_i s_{gi}}{\sum_gm_g},\qquad
 \bar s_{\mathrm{user}}=\frac1G\sum_g\frac{\sum_i s_{gi}}{m_g}.
\end{equation}
按用户重采样可以支持两种目标，但每次聚合的权重必须匹配预定目标。

多个案例、多个种子构成交叉重复测量。同一案例跨种子相关，同一种子运行内也可能共享状态。只在每个种子内部聚合再按种子重采样，回答的是固定测试集上的运行变异；若要推断新案例总体，还需按案例来源重采样，或采用适合交叉结构的分层模型与双向方案。增加重采样次数只能减少 Monte Carlo 误差，不能创造更多独立用户或独立种子。

\begin{algorithm}[配对簇重采样区间]
\textbf{输入：}两方案在同一组案例上的逐例结果、独立簇标识、目标聚合函数 $T$、重复次数 $R$、水平 $1-\alpha$。\par
\textbf{输出：}配对差点估计及百分位区间。\par
\textbf{状态：}完整的 $G$ 个簇与重采样差值列表。
\begin{enumerate}
 \item 固定缺失、超时和不适用规则；以相同案例键对齐两方案。确认簇的独立性依据及聚合权重。
 \item 计算原始差 $\hat\Delta=T(A)-T(B)$。
 \item 对 $r=1,\ldots,R$，从 $G$ 个簇有放回抽取 $G$ 次；重复抽到的簇作为重复副本加入，保留其全部配对记录。
 \item 在同一组簇副本上分别计算 $T(A^*)$ 与 $T(B^*)$，保存差值。不能独立为两个方案抽样。
 \item 取差值分布的 $\frac{\alpha}{2}$ 与 $1-\frac{\alpha}{2}$ 分位点，连同 $G$、$R$、分位数约定及点估计返回。
\end{enumerate}
\textbf{停止与边界：}完成 $R$ 次后停止；若独立簇太少或统计量在大量副本中未定义，应报告不稳定及原因，不用更多重复次数掩盖问题。百分位方法也有有限样本和偏差局限。
\end{algorithm}

\section{污染、切片及选择偏差}

训练集用于拟合，开发集用于方案选择，校准集用于阈值或裁判标定，最终测试集评估冻结方案。公开题目、答案解析、近重复、来源标题以及文件名都可能泄露信息。测试污染不仅是文本完全相同，也可能是同一模板、同一文档或答案结构进入开发过程。数据血缘与分组切分机制见\bookxref{ch:data-engineering}。

\term{切片评估}{Slice Evaluation}{}按语言、领域、长度、风险或能力结构分析子群。若各互斥切片权重为 $w_j$、得分为 $\mu_j$，总体为 $\mu=\sum_jw_j\mu_j$。假设简单任务得分0.95、困难任务0.50，简单任务占80\%时总体0.86，占50\%时总体0.725。模型在任何切片都未变化，总分却下降0.135。因此比较时需同时报告切片得分和构成变化。

重叠切片不能把样本数相加当作独立总量。事后发现的切片适合诊断，应标记探索性并在新数据中复验。高严重度失败的硬性约束需要由任务损失定义；样本量根据错误率、精度和功效要求设计。

反复测试许多方案再报告最大分数会产生选择偏差。若 $M$ 个有效检验各用水平 $\alpha$，即使每个零假设都成立，也不能把整体误报率当作 $\alpha$。Bonferroni 原则给每项 $\frac{\alpha}{M}$，利用并集上界控制总体误报概率，但可能保守。实际应预先声明主要目标、保留独立最终集，并将探索结果与确认性结论分开。未显著不等于等效；非劣性要求相对于预定损失界单独设计比较。

\section{在线实验及多目标决策}
\optional
\subsection{随机化及因果效果}
\term{随机对照实验}{Randomized Controlled Experiment}{}将合格单位随机分配到不同系统。令 $Y_i(1),Y_i(0)$ 为同一单位在两种系统下的潜在结果，目标平均处理效应为 $\mathbb E[Y(1)-Y(0)]$。随机分配使分组与潜在结果独立；在无干扰、处理定义一致和观测完整等条件下，组均值差估计该效应。

随机化单位应与记忆和干扰结构一致。同一用户跨方案可能发生学习、缓存或会话串扰；用户之间存在网络效应时，仅按用户分流也未必足够。稳定哈希可以实施固定分组，但不能替代对身份、排除规则、互斥实验和样本比例异常的处理。

例如两组各1000个独立用户，A成功600、B成功650，差值为0.05。用独立二项方差估计，标准误为
\begin{equation}
 \sqrt{\frac{0.6\times0.4}{1000}+\frac{0.65\times0.35}{1000}}\approx0.0216,
\end{equation}
近似95\%区间为 $[0.0076,0.0924]$。这与离线同题配对是不同设计，不能套用同一协方差结构。区间下界低于两个百分点，因此尚未达到该实用提升标准。

\term{意向处理分析}{Intention-to-Treat Analysis}{ITT}按最初分配组统计效果，避免只保留成功返回或完成使用的样本所造成的处理后选择偏差。超时、未完成请求和标签延迟需按预先协议处理。固定观察窗口、样本量和停止规则后再查看结果；未经校正地反复查看显著性并提前停止，会改变误报率。

在两组样本量相等、共同基率近似为 $p$ 的简化设计中，检测绝对差 $\delta$、双侧水平 $\alpha$、功效 $1-\beta$，每组样本量的粗略正态近似为
\begin{equation}
 n\approx\frac{2p(1-p)(z_{1-\frac{\alpha}{2}}+z_{1-\beta})^2}{\delta^2}.
\end{equation}
例如 $p=0.5,\delta=0.05,z_{0.975}=1.96,z_{0.8}\approx0.84$，得到每组约1568。聚类、损耗、稀有事件或多重指标会改变设计，应按实际设计调整样本量。

\subsection{质量、风险、延迟及成本}
对质量越高越好、成本与延迟越低越好的目标，若 A 各项不劣于 B 且至少一项严格更优，则 A 支配 B。\term{帕累托前沿}{Pareto Frontier}{}由所有不被支配的方案构成。构造三个方案：A为质量0.90、成本2、延迟200；B为0.92、3、240；C为0.88、2.5、220，其中成本和延迟均采用同一任意单位。A支配C，A与B互不支配。若质量下限0.91，A虽在前沿上仍不可行。

成本需包含生成、检索、评审、人工与失败重试；延迟需固定输入长度、输出长度、并发、缓存及工具路径。采用最近秩定义时，P95是排序后第 $\lceil0.95n\rceil$ 项；小样本尾部分位数不稳定，也不宜与不同实现的插值分位数混比。质量与成本均带不确定性时，还应避免仅按点估计宣布稳健支配。

\begin{note}[证据不足是一种有效决策]
发布判断结合硬约束、效应标准与统计不确定性。权限要求单独验收，罕见失败通过风险上界分析，等效或非劣性按预设界限检验。证据不足时，可维持当前方案并补充针对性评估。
\end{note}

可复现记录应绑定数据快照、模型制品、提示与索引、生成条件、评审协议、逐例输出、统计计划和环境。哈希标识制品身份，质量由任务与系统评估确定。\footnote{固定随机种子主要用于约束与重放随机过程；跨硬件、框架和并行内核时，仍可能存在数值差异，故复现实验还需声明环境与容差。}保留逐例记录才能重算指标并追溯失败，同时应控制敏感评测材料的访问与删除。上线后输入分布、标签和裁判误差持续变化，重新评估应针对实际变化确定范围。


% BEGIN GENERATED INTERVIEWS

% END GENERATED INTERVIEWS
\chapterreferences
\end{document}
